% !TEX root = ../oddp.tex

\section{The periodic and the minimal resolutions}\label{s:per_to_minimal}

\begin{warning}
Recall from Section \cref{s:simplices} and its three warnings that we are not following the usual convention to denote the simplex category.
\end{warning}

Except for the last part of this section, we assume that $r$ is odd. In \cref{prop:chains_to_minimal} we construct a chain map from $\chains(\simplex^{r})$ to $\Waug{r}$. In \cref{prop:per_to_minimal} we extend it to a chain map from $\Per{r}$ to $\Waug{r}$ that intertwines the endomorphisms $\theta$ of each of these chain complexes. Recall that $\tilde{r} = \lfloor \frac{r-1}{2}\rfloor$.

\subsection{From the simplex to the minimal resolution}

Define the following elements of $R[\Cyc_r]$ for $0\leq a<b$:
\begin{align*}
	\phi_{e}(a,b) &= \begin{cases}
		\displaystyle \sum_{i=1}^{\frac{b-a}{2}} \rho^{a+2i} & \text{if $b-a$ is even,} \\
		0 & \text{otherwise,}
	\end{cases}
	\\
	\phi_{o}(a,b) &= \begin{cases}
		\displaystyle \sum_{i=1}^{\frac{b-a-1}{2}} \rho^{a+2i+1} & \text{if $b-a$ is odd,} \\
		0 & \text{otherwise.}
	\end{cases}
\end{align*}

We say that a sequence $(a_0,\dots, a_{q-1})$ of elements in $\{0, 1, \dots, r-1\}$ is \emph{$r$-cyclically ordered} if there are integers $\bar{a}_1,\dots,\bar{a}_{q-1}$ that differ from $a_1,\dots,a_{q-1}$ by a multiple of $r$ such that 
\[
a_0=:\bar{a}_0<\bar{a}_1<\dots<\bar{a}_{q-1}<\bar{a}_q:=a_0+r.
\]
Observe that removing an element from the sequence keeps the property of being $r$-cyclically ordered.

A representative $[a_0,\dots,a_{q-1}]$ of a generator in $\chains(\partial \simplex^{r})$ is a \emph{cyclically ordered representative} if the sequence $(a_0,\dots,a_{q-1})$ is $r$-cyclically ordered. Let $A=[a_0,\dots,a_{q-1}]\in \chains(\partial \simplex^{r})$ be a cyclically ordered representative. Say that $a_i$ is \emph{even \emph{(resp. \emph{odd})} in $A$} if $\bar{a}_{i+1}-\bar{a}_i$ is an even number (resp. an odd number). 
Define the following elements of $\rW(r)$:
\begin{align*}
	\Phi_e([a_0,\dots,a_{q-1}]) &= \begin{cases}
		(-1)^{j+1}\phi_e(\bar{a}_j,\bar{a}_{j+1})\cdot e_q & \text{if $A$ has a single even entry $a_j$,} \\
		0 & \text{otherwise,}
	\end{cases}
	\\
	\Phi_o([a_0,\dots,a_{q-1}]) &= \begin{cases}
		\sum_{j=0}^{q-1} \phi_o(\bar{a}_j,\bar{a}_{j+1})\cdot e_q & \text{if $A$ has no even entries,} \\
		0 & \text{otherwise.}
	\end{cases}
\end{align*} 
In what follows we work always with cyclic representatives. Let $\varphi(q) = \lfloor\frac{r-q-1}{2}\rfloor$ if $q<r-1$ and $\varphi(r-1) =1$.

\begin{proposition}\label{prop:chains_to_minimal} The linear homomorphism with $\bZ[\frac{1}{\tilde{r}!}]$-coefficients
	\begin{align*}
		\Phi \colon \chains(\partial \simplex^{r})& \lra \Waug{r}
		&
		\Phi(A) = \begin{cases}
			e_0 & \text{if $q=0$ and $A = \emptyset$,}\\
			\frac{\varphi(q)!}{\tilde{r}!}\Phi_e(A) & \text{if $q$ is even and positive,} \\
			\frac{\varphi(q)!}{\tilde{r}!}\Phi_o(A) & \text{if $q$ is odd.}
		\end{cases}
	\end{align*}
	is a chain map.
\end{proposition}

\begin{remark}
	If $r=3$, then $\bZ[\frac{1}{\tilde{r}!}] = \bZ$ and if $r$ is prime, then $\mathbb{F}_r$ is a $\bZ[\frac{1}{\tilde{r}!}]$-algebra, hence these coefficients are also allowed.
\end{remark}

\begin{proof} In order to simplify the notation, we will omit the generator $e_q$ in the formulas below.
	Suppose first that $A = [a_0,\dots,a_{q-1}]$ has even length, in which case $\varphi(q) = \varphi(q-1)$. Since $A$ has even length, it must have at least one even entry. If $A$ has more than one even entry, then all sequences in its boundary have at least one even entry, and therefore $\Phi_o(\partial A) = 0 = \partial\Phi_e(A)$. If $A$ has a single even entry $a_j$, then
	\begin{align*}
		\partial \Phi_e(A) &= (-1)^{j+1}\partial \phi_e(\bar a_j,\bar a_{j+1}) = (-1)^{j+1}\left(\rho\phi_e(\bar a_j,\bar a_{j+1}
		) - \phi_e(\bar a_j,\bar a_{j+1})\right).
	\end{align*}
	On the other hand, since $[a_0,\dots,\hat{a}_i,\dots,a_{q-1}]$ will have two even entries unless $i=j$ or $i=j+1$, we have that
	\begin{align*}
		\Phi_o(\partial A)
		&= (-1)^{j}\left(\Phi_o ([a_0,\dots,\hat{a}_j,\dots,a_{q-1}])-\Phi_o ([a_0,\dots,\hat{a}_{j+1},\dots,a_{q-1}])\right)
	\end{align*}
	Now, using that $\phi_o(\bar a_{j-1},\bar a_{j+1}) = \phi_o(\bar a_{j-1},\bar a_j) + \phi_e(\bar a_j,\bar a_{j+1})$, the first summand equals
	\[
	\phi_o(\bar a_{j-1},\bar a_j) + \phi_e(\bar a_j,\bar a_{j+1})+\sum_{i\neq j-1,j} \phi_o(\bar a_i,\bar a_{i+1}).
	\]
	Using that $\phi_o(\bar a_{j},\bar a_{j+2}) = \phi_e(\bar a_j+1,\bar a_{j+1}+1)+\phi_o(\bar a_{j+1},\bar a_{j+2})$, the second summand equals
	\[
	\phi_e(\bar a_j+1,\bar a_{j+1}+1) + \phi_o(\bar a_{j+1},\bar a_{j+2})+\sum_{i\neq j,j+1} \phi_o(\bar a_i,\bar a_{i+1}).
	\]
	Therefore, most terms cancel and we are left with the value of $\partial \Phi_e([a_0,\dots,a_{q-1}])$.

	Suppose now that $A = [a_0,\dots,a_{q-1}]$ has odd length, in which case $\tilde{r}-\tilde{q}-1 = \tilde{r}-\widetilde{q-1} $. If $A$ has more than two even entries, then $[a_0,\dots,\hat{a}_i,\dots,a_{q-1}]$ has more than one even entry, hence $\Phi_e([a_0,\dots,\hat{a}_i,\dots,a_{q-1}]) = 0$ for all $i$, therefore $\Phi_e(\partial A) = 0 = \partial \Phi_o(A)$. If $A$ has exactly two even entries $a_j$, $a_k$, then we have two cases: if these entries are not consecutive and $j<k$, then the nontrivial summands in $ \Phi_e(\partial A)$ are
	\begin{align*}
		(-1)^j\Phi_e([a_0,\dots,\hat{a}_j,\dots,a_{q-1}]) &= (-1)^j(-1)^{k}\phi_e(\bar a_k,\bar a_{k+1}) \\
		(-1)^{j+1}\Phi_e([a_0,\dots,\hat{a}_{j+1},\dots,a_{q-1}]) &= (-1)^{j+1}(-1)^{k}\phi_e(\bar a_k,\bar a_{k+1})
		\\
		(-1)^{k}\Phi_e([a_0,\dots,\hat{a}_k,\dots,a_{q-1}]) &= (-1)^k(-1)^{j+1}\phi_e(\bar a_j,\bar a_{j+1}) \\
		(-1)^{k+1}\Phi_e([a_0,\dots,\hat{a}_{k+1},\dots,a_{q-1}]) &= (-1)^{k+1}(-1)^{j+1}\phi_e(\bar a_j,\bar a_{j+1})
	\end{align*}
	which cancel in pairs. If the entries are consecutive, then $k=j+1$ and the nontrivial summands in $\Phi_e(\partial A)$ are
	\begin{align*}
		(-1)^j\Phi_e([a_0,\dots,\hat{a}_j,\dots,a_{q-1}]) &= (-1)^{j}(-1)^{j+1}\phi_e(\bar a_{j+1},\bar a_{j+2}) \\
		(-1)^{j+1}\Phi_e([a_0,\dots,\hat{a}_{j+1},\dots,a_{q-1}]) &= (-1)^{j+1}(-1)^{j+1}\phi_e(\bar a_j,\bar a_{j+2}) \\
		(-1)^{j+2}\Phi_e([a_0,\dots,\hat{a}_{j+2},\dots,a_{q-1}]= &= (-1)^{j+2}(-1)^{j+1}\phi_e(\bar a_j,\bar a_{j+1})
	\end{align*}
	and the middle summand cancels with the other two, so $\Phi_e(\partial (A)) = 0 = \partial(\Phi_o(A))$.

	If $A$ has no even entry, then
	\[
	\partial \Phi_o(A) = N\left(\sum_{i=0}^{q-1} \phi_o(\bar a_i,\bar a_{i+1})\right) = \frac{r-q}{2}N
	\]
	On the other hand,
	\[
	\Phi_e(\partial A)
	= \sum_{i=0}^{q-1} \phi_e(\bar a_i,\bar a_{i+2})
	= N
	\]
	This completes the proof, since $\frac{r-q}{2} = \frac{\varphi(q-1)!}{\varphi(q)!}$.
\end{proof}

\begin{example}\label{example:Phi}
If $r=3$, we have:
\begin{align*}
\Phi([0]) &= e_1
&
\Phi([1,2]) &= \rho e_2
\end{align*}
If $r=5$, we have:
\begin{align*}
\Phi([0,1,2]) &= 2^{-1}e_3
&
\Phi([0,1]) &= 2^{-1}(e_2 + \rho^3e_2)
&
\Phi([0,3]) &= 2^{-1}e_2.
\end{align*}
\end{example}

\begin{remark}\label{remark:phidual}
	For each $0\leq q<r-1$, let $L_q \subset (\partial \simplex^{r})_q$ be the set of those simplices $[a_0,\dots, a_{q-1}]$ with increasing entries (in the total order) such that $a_{2i}$ is even and different from $r-1$, and $a_{2i+1}$ is odd. Then the linear dual of $\Phi$ has the following expression
	\begin{equation}\label{eq:111}
		\Phi^\vee(e_{q}^\dd) = \frac{\varphi(q)!}{\tilde{r}!}\sum_{A\in L_q} A^\vee.
	\end{equation}
\end{remark}

\begin{proposition}\label{prop:per_to_minimal}
	The chain map $\Phi$ extends to a chain map $\Phi\colon \Per{r}\to \Waug{r}$ defining $\Phi(\sus{(r-1)m}\tau) = (\tilde{r}!)^{-m}\theta^{m}\Phi(\tau)$. As a consequence, the map 
	\[
		\phi\colon \sus{\bullet (r-1)}\Per{r}\lra \sus{\bullet(r-1)}\Waug{r}
	\] 
	given by $\phi_n(\tau) = (\tilde{r}!)^n\Phi(\tau)$ is a cosemi-simplicial map well-defined with any coefficients.
\end{proposition}
\begin{proof}                                                                          
	Let $v = [i]$ be a vertex of $\simplex^{r}$. We need to prove that $\Phi\circ \partial(\sus{r-1}v) = \partial\circ \Phi(\sus{r-1}(v))$. The right hand-side is 
	\begin{align*}
		\frac{1}{\tilde{r}!} \partial\circ \theta \circ \Phi(v) &= 
		\frac{1}{\tilde{r}!} \partial\circ \theta \left(\frac{1}{\tilde{r}}\phi_o(i,i+r)\right) \\
		&= \frac{1}{\tilde{r}!} \partial\circ \theta \left(\frac{1}{\tilde{r}}\sum_{j=1}^{\frac{r-1}{2}}\rho^{i+2j+1} e_{1}\right) \\
		&= \frac{1}{\tilde{r}!} \theta\circ\partial \left(\frac{1}{\tilde{r}}\sum_{j=1}^{\frac{r-1}{2}}\rho^{i+2j+1} e_{1}\right) \\
		&= \frac{1}{\tilde{r}!} \theta\circ N(e_0) \\
		&= \frac{1}{\tilde{r}!} N(e_{r-1}) \\
		%&= \frac{1}{\tilde{r}!} \circ \theta (\sum_{i=1}^{\frac{r-1}{2}})\rho^{j+2i+1} \\
		\end{align*}
	The left hand-side is
	\begin{align*}
		\Phi(\partial\sigma) &= \sum_{j=0}^{r-1} (-1)^j\Phi([0,\ldots,\hat{j},\ldots,r-1]) \\
		&= \frac{1}{\tilde{r}!}\left((-1)^0(-1)^{r-1}\phi_e(r-1,r+1)+\sum_{j=1}^{r-1} (-1)^j(-1)^{j}\phi_e(j-1,j+1)\right)e_{r-1} \\
		&= \frac{1}{\tilde{r}!}\left (\rho^{r+1}+\sum_{j=1}^{r-1} \rho^{j+1}\right) e_{r-1} \\
		&= \frac{1}{\tilde{r}!}N(e_{r-1}).
	\end{align*}
	Finally, we check that $\phi$ commutes with the coface maps $\d^i$ (recall that these do not depend on $i$): if $\tau\neq \varnothing$,
	 \begin{align*}
	\phi\circ \d^i(\tau) &= (\tilde{r}!)^{n+1}\Phi(\sus{r-1}\tau) \\
	&= (\tilde{r}!)^{n+1}\cdot (\tilde{r}!)^{-1}\theta\Phi(\tau)\\
	&= \theta\left((\tilde{r}!)^{n}\Phi(\tau)\right).\\
	&= \d^i\circ \phi(\tau).
	\end{align*}
	If $\tau=\emptyset$, we have that $\d^i(\tau) = \sum(-1)^j [0,1,\dots,\hat{j},\ldots,r]$ and one can check from the definition of $\phi$ that $\phi([0,1,\dots,\hat{i},\ldots,r]) = e_{r-1}$.
\end{proof}

